
% --- SLIDE 3: HIDDEN TECHNICAL DEBT ---
\begin{frame}{Hidden Technical Debt}

    \begin{block}{The core argument (Sculley et al., NeurIPS 2015)}
        Building an ML system is fast and cheap. \textbf{Maintaining} it is slow and expensive.
        The debt is \emph{hidden} because it lives at the system level, not in the code, so ordinary
        code review and unit tests do not find it.
    \end{block}

    \vspace{0.5em}
    \begin{columns}[T]
        \begin{column}{0.48\textwidth}
            \small
            \textbf{Boundary erosion}
            \begin{itemize}
                \item \textbf{Entanglement (CACE):} \textit{Changing Anything Changes Everything}. Change one
                      feature's distribution and every other weight shifts.
                \item \textbf{Correction cascades:} a model that patches another model's output --- now you cannot
                      improve either one alone.
                \item \textbf{Undeclared consumers:} someone reads your predictions off a log and you never find out.
            \end{itemize}
        \end{column}
        \begin{column}{0.48\textwidth}
            \small
            \textbf{System-level anti-patterns}
            \begin{itemize}
                \item \textbf{Glue code:} 95\% of the repository is plumbing around a 5\% library call.
                \item \textbf{Pipeline jungles:} scrapes and joins accreted over years; nobody can rerun them.
                \item \textbf{Configuration debt:} a hundred flags, no schema, no review.
                \item \textbf{Feedback loops:} the model influences the data it will next be trained on.
            \end{itemize}
        \end{column}
    \end{columns}

\end{frame}
